\documentclass[12pt]{article}

% Match the earlier reports using locally available fonts.
\usepackage[T1]{fontenc}
\usepackage{times}
\usepackage{microtype}
\usepackage{amsmath}
\usepackage[parfill]{parskip}
\setlength{\parindent}{0pt}
\setlength{\parskip}{2ex plus 0.2ex minus 0.2ex}
\usepackage[letterpaper,margin=1in]{geometry}
\usepackage[hidelinks]{hyperref}

\begin{document}

\begin{flushleft}
\textbf{Week 4: Markov Chain Monte Carlo} \\
Jaisel Singh (js6897) \\
October 7, 2026
\end{flushleft}

\vspace{0.1in}

What stood out to me in Neal's review was that Bayesian inference often comes
down to averaging over plausible parameter values. This changes the goal from
last week's optimization: a single high-probability point does not describe the
uncertainty around it. MCMC constructs a Markov chain whose stationary
distribution is the target posterior, allowing us to estimate expectations from
its draws when the chain explores adequately. I found it especially useful that
Metropolis acceptance ratios cancel the unknown normalizing constant. We can
therefore work with an unnormalized posterior even when computing its integral
is impractical.

The gap between a correct algorithm and a useful finite run also stood out.
Successive draws are dependent, and strong positive autocorrelation means that
collecting many draws may add little information. I found effective sample size
a more meaningful way to judge precision than the raw iteration count. A small
Monte Carlo standard error does not mean the posterior itself is narrow:
one measures uncertainty in our numerical estimate, while the other reflects
uncertainty about the model's parameters. This distinction makes me more careful
about what it means to say that a result is precise.

Hoffman and Gelman's No-U-Turn Sampler addresses some of this computational
difficulty. Hamiltonian Monte Carlo uses momentum and gradients to propose
moves that can travel farther through the distribution than a random walk.
NUTS chooses trajectory lengths automatically by stopping expansion when a path
starts doubling back, while preserving the target distribution through its
candidate-selection rules. I like how this makes a sophisticated sampler easier
to use, though automatic tuning does not establish convergence. My remaining
question is how to distinguish a chain that has explored the posterior well
from one that looks stable because it is stuck near a single mode.

\vspace{0.05in}
{\footnotesize
\textit{Readings:} Neal (1993), \href{https://ftp.cs.toronto.edu/pub/radford/review.pdf}{\emph{Probabilistic Inference Using Markov Chain Monte Carlo Methods}};
Hoffman and Gelman (2014), \href{https://jmlr.org/papers/v15/hoffman14a.html}{\emph{The No-U-Turn Sampler: Adaptively Setting Path Lengths in Hamiltonian Monte Carlo}}.
\par}

\end{document}

%%% Local Variables:
%%% mode: latex
%%% TeX-master: t
%%% End:
